\chapter{数学基础：从概率到李群}\label{ch:02-math}

\chapterintro{这一章不追求数学上的完备，只准备“够用”的工具箱。我们按定位与导航里真正会用到的东西来组织：坐标系怎么描述、概率怎么描述不确定性、方差和协方差到底在算什么、如何用最小二乘去拟合一个位姿、以及为什么旋转不能直接相加而要用李群。每一个概念都会从定义讲起，给出计算公式和一个可以手算的小例子。读完这一章，后面所有章节里的公式你至少能“看懂每一行在做什么”。}

\section{坐标系与刚体变换}\label{sec:02-frame}

\subsection{点、向量与坐标系}

一个三维空间中的点，只有在指定了坐标系之后才有坐标。本书约定：用 $\{W\}$ 表示世界系（map 系），$\{B\}$ 表示机体系（默认与 IMU 系重合），$\{L\}$、$\{C\}$ 分别表示雷达系和相机系。点 $P$ 在 $\{B\}$ 系下的坐标记作 ${}^{B}\bm{p}\in\R^{3}$，写成列向量
\begin{equation}
{}^{B}\bm{p}=\begin{bmatrix}p_x\\p_y\\p_z\end{bmatrix}.
\end{equation}
同一个点在不同坐标系下的坐标不同，它们之间通过“旋转 + 平移”联系起来：
\begin{equation}
{}^{W}\bm{p}={}^{W}\bm{R}_{B}\,{}^{B}\bm{p}+{}^{W}\bm{p}_{B},
\label{eq:02-rigid}
\end{equation}
其中 ${}^{W}\bm{R}_{B}\in\R^{3\times3}$ 是旋转矩阵，${}^{W}\bm{p}_{B}\in\R^{3}$ 是 $B$ 系原点在 $W$ 系中的位置。式~\eqref{eq:02-rigid} 就是刚体变换（rigid transformation）的全部内容：先转，再平移。

\subsection{旋转矩阵}

\begin{keybox}{旋转矩阵的定义}
一个 $3\times3$ 矩阵 $\bm{R}$ 是旋转矩阵，当且仅当它满足
\begin{equation}
\bm{R}^{\top}\bm{R}=\bm{I},\qquad \det(\bm{R})=1.
\label{eq:02-so3def}
\end{equation}
所有这样的矩阵构成特殊正交群 $\SO(3)$。
\end{keybox}

为什么是这两个条件？$\bm{R}^{\top}\bm{R}=\bm{I}$ 保证长度和夹角不变：对任意向量 $\bm{v}$，有 $\norm{\bm{R}\bm{v}}^{2}=\bm{v}^{\top}\bm{R}^{\top}\bm{R}\bm{v}=\bm{v}^{\top}\bm{v}=\norm{\bm{v}}^{2}$。而 $\det(\bm{R})=\pm1$ 中取 $+1$ 是为了排除“镜像”，因为镜子里的世界不是旋转出来的。

旋转矩阵的复合就是矩阵乘法：先转 ${}^{B}\bm{R}_{C}$ 再转 ${}^{W}\bm{R}_{B}$，总旋转是 ${}^{W}\bm{R}_{B}{}^{B}\bm{R}_{C}$。注意下标“相邻相消”：
\begin{equation}
{}^{W}\bm{p}={}^{W}\bm{R}_{B}\left({}^{B}\bm{R}_{C}{}^{C}\bm{p}\right)
=\left({}^{W}\bm{R}_{B}{}^{B}\bm{R}_{C}\right){}^{C}\bm{p}.
\end{equation}
这个“下标相消”规则在后面的 TF 树、标定、外参里会反复出现，值得记牢。

\subsection{齐次坐标与 $\SE(3)$}

式~\eqref{eq:02-rigid} 是“先乘再加”，不便于连续复合。引入齐次坐标可以把旋转和平移合并成一次矩阵乘法。定义
\begin{equation}
{}^{W}\bar{\bm{p}}=\begin{bmatrix}{}^{W}\bm{p}\\1\end{bmatrix}\in\R^{4},\qquad
{}^{W}\bm{T}_{B}=\begin{bmatrix}{}^{W}\bm{R}_{B}&{}^{W}\bm{p}_{B}\\\bm{0}^{\top}&1\end{bmatrix}\in\R^{4\times4},
\end{equation}
于是
\begin{equation}
{}^{W}\bar{\bm{p}}={}^{W}\bm{T}_{B}\,{}^{B}\bar{\bm{p}}.
\end{equation}
所有这样的 $\bm{T}$ 构成特殊欧氏群 $\SE(3)$。复合规则同样是矩阵乘法，且
\begin{equation}
\bm{T}^{-1}=\begin{bmatrix}\bm{R}^{\top}&-\bm{R}^{\top}\bm{t}\\\bm{0}^{\top}&1\end{bmatrix}.
\label{eq:02-se3inv}
\end{equation}
验证式~\eqref{eq:02-se3inv}：直接乘开，左上角 $\bm{R}^{\top}\bm{R}=\bm{I}$，右上角 $\bm{R}^{\top}\bm{t}-\bm{R}^{\top}\bm{t}=\bm{0}$。这个逆的闭式表达比数值求逆快得多，在代码里几乎总用这个形式。

\subsection{欧拉角与万向锁}

把旋转拆成绕 $z$、$y$、$x$ 三个轴依次转的角度 $(\psi,\theta,\phi)$（yaw、pitch、roll），叫欧拉角。它直观、人类可读（IMU 输出的姿态常用它显示），但有两个致命问题：

\begin{enumerate}
  \item \textbf{万向锁（gimbal lock）}：当 pitch 接近 $\pm90^\circ$ 时，yaw 和 roll 退化成同一个自由度，姿态表示不唯一，运动学方程出现奇异。
  \item \textbf{插值困难}：两组欧拉角之间直接线性插值，得到的中间姿态往往不是“最短路径”旋转。
\end{enumerate}

因此，内部计算一律用旋转矩阵或四元数，欧拉角只在“给人看”的时候使用。

\section{概率论最小工具箱}\label{sec:02-prob}

定位的本质是“在不确定中估计”。传感器有噪声，运动模型有误差，所以我们用概率分布描述“我认为机器人在哪”。这一节把后面会用到的概率概念逐个讲清楚，尤其是\textbf{方差}和\textbf{协方差}——它们是全书出现频率最高的两个词。

\subsection{随机变量与概率}

一个随机变量 $X$ 是“取值不确定的量”。如果 $X$ 只能取有限个离散值 $x_1,x_2,\dots$，用概率质量函数 $p(X=x_i)$ 描述，满足
\begin{equation}
p(X=x_i)\ge 0,\qquad \sum_i p(X=x_i)=1.
\end{equation}
如果 $X$ 可以在实数轴上连续取值，用概率密度函数 $p(x)$ 描述，满足
\begin{equation}
p(x)\ge 0,\qquad \int_{-\infty}^{+\infty}p(x)\diff x=1,
\end{equation}
并且落在区间 $[a,b]$ 的概率是
\begin{equation}
P(a\le X\le b)=\int_a^b p(x)\diff x.
\end{equation}
注意：$p(x)$ 本身不是概率，它可能大于 1；只有积分出来的面积才是概率。

\subsection{期望}

\begin{keybox}{定义：期望}
离散随机变量的期望是
\begin{equation}
\E[X]=\sum_i x_i\,p(X=x_i),
\end{equation}
连续随机变量的期望是
\begin{equation}
\E[X]=\int_{-\infty}^{+\infty}x\,p(x)\diff x.
\label{eq:02-expect}
\end{equation}
期望就是“加权平均”，权重是概率。
\end{keybox}

\textbf{例}：掷一枚均匀硬币，正面记 1、反面记 0，则
\begin{equation}
\E[X]=1\times 0.5+0\times 0.5=0.5.
\end{equation}

期望最重要的性质是\textbf{线性性}，它对任意随机变量（无论是否独立）都成立：
\begin{equation}
\E[aX+b]=a\E[X]+b,\qquad \E[X+Y]=\E[X]+\E[Y].
\label{eq:02-lin}
\end{equation}
证明第一个式子（离散情形）：
\begin{equation}
\E[aX+b]=\sum_i (ax_i+b)p(x_i)=a\sum_i x_ip(x_i)+b\sum_i p(x_i)=a\E[X]+b.
\end{equation}
最后一步用了 $\sum_i p(x_i)=1$。这个性质是后面所有“误差传播”公式的基础。

\subsection{方差与标准差}

\begin{keybox}{定义：方差}
方差衡量随机变量偏离其期望的程度，定义为
\begin{equation}
\Var(X)=\E\left[(X-\E[X])^{2}\right].
\label{eq:02-vardef}
\end{equation}
它的平方根 $\sigma_X=\sqrt{\Var(X)}$ 叫标准差。
\end{keybox}

把定义展开，可以得到更常用的计算公式：
\begin{equation}
\Var(X)=\E\left[X^{2}-2X\E[X]+(\E[X])^{2}\right]
=\E[X^{2}]-2\E[X]\E[X]+(\E[X])^{2}
=\E[X^{2}]-(\E[X])^{2}.
\label{eq:02-var-expand}
\end{equation}
第二步用了期望的线性性，第三步用了 $\E[\E[X]]=\E[X]$（$\E[X]$ 是一个常数）。所以记住这个等价形式：
\begin{equation}
\boxed{\Var(X)=\E[X^{2}]-\left(\E[X]\right)^{2}}
\end{equation}
即“平方的期望减去期望的平方”。

\textbf{例}：掷一枚均匀骰子，$X\in\{1,2,3,4,5,6\}$，每个概率 $1/6$。
\begin{align}
\E[X]&=\frac{1+2+3+4+5+6}{6}=3.5,\\
\E[X^{2}]&=\frac{1+4+9+16+25+36}{6}=\frac{91}{6}\approx 15.17,\\
\Var(X)&=\frac{91}{6}-3.5^{2}=15.17-12.25\approx 2.92,\\
\sigma_X&=\sqrt{2.92}\approx 1.71.
\end{align}
也就是说，掷骰子的结果平均偏离 3.5 约 1.71。方差是“偏离程度的平方的平均”，开根号后才和原数据同量纲。

方差有两个常用性质，后面推导漂移时会直接用到：
\begin{equation}
\Var(aX+b)=a^{2}\Var(X),
\label{eq:02-var-scale}
\end{equation}
\begin{equation}
\text{若 }X,Y\text{ 相互独立，则}\ \Var(X+Y)=\Var(X)+\Var(Y).
\label{eq:02-var-sum}
\end{equation}
先证式~\eqref{eq:02-var-scale}：
\begin{equation}
\Var(aX+b)=\E\left[(aX+b-\E[aX+b])^{2}\right]
=\E\left[(aX-a\E[X])^{2}\right]
=a^{2}\E\left[(X-\E[X])^{2}\right]=a^{2}\Var(X).
\end{equation}
再证式~\eqref{eq:02-var-sum}。先算
\begin{equation}
\Var(X+Y)=\E\left[(X+Y)^{2}\right]-\left(\E[X+Y]\right)^{2}
=\E[X^{2}]+2\E[XY]+\E[Y^{2}]-\left(\E[X]+\E[Y]\right)^{2}.
\end{equation}
展开右边最后一项：
\begin{equation}
\left(\E[X]+\E[Y]\right)^{2}=(\E[X])^{2}+2\E[X]\E[Y]+(\E[Y])^{2}.
\end{equation}
两式相减得
\begin{equation}
\Var(X+Y)=\left[\E[X^{2}]-(\E[X])^{2}\right]+\left[\E[Y^{2}]-(\E[Y])^{2}\right]+2\left(\E[XY]-\E[X]\E[Y]\right).
\end{equation}
前两个方括号分别是 $\Var(X)$、$\Var(Y)$；而独立意味着 $\E[XY]=\E[X]\E[Y]$，最后一项为零。于是 $\Var(X+Y)=\Var(X)+\Var(Y)$。

\subsection{协方差与协方差矩阵}

方差只描述一个变量。两个变量“一起变化”的程度用协方差描述：
\begin{equation}
\Cov(X,Y)=\E\left[(X-\E[X])(Y-\E[Y])\right]=\E[XY]-\E[X]\E[Y].
\label{eq:02-cov}
\end{equation}
展开方式与方差完全一样。$\Cov(X,X)=\Var(X)$。若 $X,Y$ 独立，则 $\Cov(X,Y)=0$；但反过来不成立——协方差为零只说明“没有线性相关”，可能存在非线性关系。

对多维随机向量 $\bm{x}=[x_1,\dots,x_n]^{\top}$，把所有两两协方差排成矩阵：
\begin{equation}
\bm{\Sigma}=\E\left[(\bm{x}-\E[\bm{x}])(\bm{x}-\E[\bm{x}])^{\top}\right]\in\R^{n\times n},
\label{eq:02-covmat}
\end{equation}
叫协方差矩阵。它的对角元是各方差 $\Sigma_{ii}=\Var(x_i)$，非对角元是协方差 $\Sigma_{ij}=\Cov(x_i,x_j)$。协方差矩阵一定是对称半正定的，因为对任意向量 $\bm{a}$，
\begin{equation}
\bm{a}^{\top}\bm{\Sigma}\bm{a}=\E\left[\left(\bm{a}^{\top}(\bm{x}-\E[\bm{x}])\right)^{2}\right]\ge 0.
\end{equation}

\textbf{协方差的几何意义}：二维情况下，协方差矩阵的特征向量给出分布的“主轴方向”，特征值给出沿主轴的方差大小。特征值越大，说明该方向上不确定性越大，分布椭圆越长。第 9 章用“马氏距离”做外点剔除、第 13 章用“信息矩阵”加权残差，本质都是在利用这个椭圆。

\subsection{高斯分布}

一维高斯（正态）分布的概率密度为
\begin{equation}
p(x)=\frac{1}{\sqrt{2\pi}\,\sigma}\exp\left(-\frac{(x-\mu)^{2}}{2\sigma^{2}}\right),
\label{eq:02-gauss1d}
\end{equation}
记作 $x\sim\mathcal{N}(\mu,\sigma^{2})$。它由两个参数完全确定：均值 $\mu$ 和方差 $\sigma^{2}$。

\textbf{验证它确实是概率密度}：首先换元 $u=(x-\mu)/(\sqrt{2}\sigma)$，则 $\diff x=\sqrt{2}\sigma\diff u$，
\begin{equation}
\int_{-\infty}^{+\infty}p(x)\diff x=\frac{1}{\sqrt{2\pi}\sigma}\int_{-\infty}^{+\infty}e^{-u^{2}}\sqrt{2}\sigma\diff u
=\frac{1}{\sqrt{\pi}}\int_{-\infty}^{+\infty}e^{-u^{2}}\diff u.
\end{equation}
利用高斯积分 $\int_{-\infty}^{+\infty}e^{-u^{2}}\diff u=\sqrt{\pi}$，得到积分等于 1。其次算期望：
\begin{equation}
\E[X]=\int x\,p(x)\diff x=\mu,
\end{equation}
因为密度关于 $x=\mu$ 对称。再算方差，令 $u=(x-\mu)/\sigma$：
\begin{equation}
\Var(X)=\frac{1}{\sqrt{2\pi}\sigma}\int_{-\infty}^{+\infty}(x-\mu)^{2}e^{-\frac{(x-\mu)^{2}}{2\sigma^{2}}}\diff x
=\frac{\sigma^{2}}{\sqrt{2\pi}}\int_{-\infty}^{+\infty}u^{2}e^{-u^{2}/2}\diff u=\sigma^{2},
\end{equation}
其中用到分部积分得到的 $\int u^{2}e^{-u^{2}/2}\diff u=\sqrt{2\pi}$。

多维高斯分布写成
\begin{equation}
p(\bm{x})=\frac{1}{(2\pi)^{n/2}\abs{\bm{\Sigma}}^{1/2}}
\exp\left(-\frac{1}{2}(\bm{x}-\bm{\mu})^{\top}\bm{\Sigma}^{-1}(\bm{x}-\bm{\mu})\right),
\label{eq:02-gaussnd}
\end{equation}
记作 $\bm{x}\sim\mathcal{N}(\bm{\mu},\bm{\Sigma})$。指数里的二次型
\begin{equation}
d^{2}=(\bm{x}-\bm{\mu})^{\top}\bm{\Sigma}^{-1}(\bm{x}-\bm{\mu})
\label{eq:02-mahal}
\end{equation}
叫马氏距离的平方。它和欧氏距离的区别是“按不确定性的方向加权”：方差小的方向权重更大。第 9 章会证明，$d^{2}$ 服从卡方分布，这正是外点剔除检验的依据。

\subsection{条件概率与贝叶斯公式}

条件概率的定义是
\begin{equation}
p(A\mid B)=\frac{p(A,B)}{p(B)},\qquad p(B)>0,
\label{eq:02-cond}
\end{equation}
即“在 $B$ 发生的条件下 $A$ 发生的概率”等于联合概率除以 $B$ 的概率。由对称性 $p(A,B)=p(B\mid A)p(A)$，代入式~\eqref{eq:02-cond} 并交换 $A$、$B$ 的位置，得到贝叶斯公式：
\begin{equation}
p(A\mid B)=\frac{p(B\mid A)\,p(A)}{p(B)}.
\label{eq:02-bayes}
\end{equation}
把它放到定位里：$A$ 是“机器人的位姿”，$B$ 是“传感器观测”。于是
\begin{equation}
\underbrace{p(\text{位姿}\mid\text{观测})}_{\text{后验}}
=\frac{\underbrace{p(\text{观测}\mid\text{位姿})}_{\text{似然}}\underbrace{p(\text{位姿})}_{\text{先验}}}
{\underbrace{p(\text{观测})}_{\text{证据}}}.
\end{equation}
后验就是“看到这帧点云后，我认为自己在哪”。分母与位姿无关，在优化里通常被忽略（当作常数）。整本书的定位与 SLAM 算法，本质上都在算这个后验：滤波器用递推近似它，优化器用最小二乘找它的最大值。

\subsection{独立性}

两个随机变量独立，当且仅当
\begin{equation}
p(X,Y)=p(X)p(Y)\quad\Longleftrightarrow\quad p(X\mid Y)=p(X).
\end{equation}
独立是后面大量公式成立的前提：多个观测独立时，似然可以连乘，取对数后变成求和，最小二乘的各项才能分开写。工程上要特别小心——时间上相邻的传感器读数、同一平面上相邻的点，往往并不独立，第 23 章 Z-Localization 讨论的“平面内点残差相关”就是一例。

\subsection{卡方分布}

\begin{definition}[卡方分布]
设 $X_1,X_2,\dots,X_k$ 是相互独立的标准正态随机变量，$X_i\sim\mathcal{N}(0,1)$，则
\begin{equation}
Q=X_1^{2}+X_2^{2}+\cdots+X_k^{2}
\end{equation}
服从自由度为 $k$ 的卡方分布，记作 $Q\sim\chi^{2}(k)$。
\end{definition}

自由度 $k=1$ 时密度函数为
\begin{equation}
p(q)=\frac{1}{\sqrt{2\pi q}}e^{-q/2},\qquad q\ge 0,
\end{equation}
$k=2$ 时是 $p(q)=\frac12 e^{-q/2}$，$k$ 一般时
\begin{equation}
p(q)=\frac{1}{2^{k/2}\Gamma(k/2)}q^{k/2-1}e^{-q/2},
\end{equation}
其中 $\Gamma$ 是伽马函数（$\Gamma(n)=(n-1)!$）。我们不要求记住这个密度，只需要知道\textbf{给定置信水平 $\alpha$ 时，分位数 $\chi^{2}_{k,\alpha}$ 可以从表里查到}。

为什么工程里关心卡方分布？因为如果残差 $\bm{r}$ 满足零均值高斯分布、协方差为 $\bm{S}$，那么二次型
\begin{equation}
d^{2}=\bm{r}^{\top}\bm{S}^{-1}\bm{r}
\end{equation}
就是若干个标准正态变量的平方和，服从 $\chi^{2}(k)$（$k$ 是残差维数）。于是可以设定一个阈值 $\chi^{2}_{k,0.95}$：如果 $d^{2}$ 超过它，说明这个观测“太离谱”，应当作为外点剔除。这就是第 9 章马氏距离检验、第 23 章 NIS 检验的数学基础。

\section{微积分与优化}\label{sec:02-opt}

\subsection{泰勒展开}

一元函数在 $x_0$ 处的泰勒展开是
\begin{equation}
f(x_0+\delta)=f(x_0)+f'(x_0)\delta+\frac{1}{2}f''(x_0)\delta^{2}+\frac{1}{6}f'''(x_0)\delta^{3}+\cdots
\label{eq:02-taylor}
\end{equation}
它的含义是：用多项式在局部逼近函数。当我们只关心很小的 $\delta$ 时，可以只保留前两项（一阶近似）或前三项（二阶近似），把高阶项写成 $O(\delta^{2})$、$O(\delta^{3})$。

多元函数的一阶近似为
\begin{equation}
f(\bm{x}_0+\delta\bm{x})\approx f(\bm{x}_0)+\nabla f(\bm{x}_0)^{\top}\delta\bm{x},
\end{equation}
其中梯度 $\nabla f=\left[\partial f/\partial x_1,\dots,\partial f/\partial x_n\right]^{\top}$。二阶近似为
\begin{equation}
f(\bm{x}_0+\delta\bm{x})\approx f(\bm{x}_0)+\nabla f^{\top}\delta\bm{x}+\frac12\delta\bm{x}^{\top}\bm{H}\delta\bm{x},
\end{equation}
其中 $\bm{H}$ 是海森矩阵（见下节）。\textbf{整本书的非线性优化，本质上就是反复做这个一阶/二阶近似，然后解一个线性问题。}

\subsection{雅可比矩阵}

\begin{definition}[雅可比矩阵]
设 $\bm{f}:\R^{n}\to\R^{m}$ 是一个向量函数，$\bm{f}(\bm{x})=[f_1(\bm{x}),\dots,f_m(\bm{x})]^{\top}$。它的雅可比矩阵 $\bm{J}\in\R^{m\times n}$ 定义为
\begin{equation}
\bm{J}=\frac{\partial\bm{f}}{\partial\bm{x}}=
\begin{bmatrix}
\dfrac{\partial f_1}{\partial x_1}&\cdots&\dfrac{\partial f_1}{\partial x_n}\\
\vdots&\ddots&\vdots\\
\dfrac{\partial f_m}{\partial x_1}&\cdots&\dfrac{\partial f_m}{\partial x_n}
\end{bmatrix}.
\end{equation}
\end{definition}
雅可比就是“多元函数的导数”，它告诉你：输入变化一点点，输出会怎么变。$\bm{f}(\bm{x}+\delta\bm{x})\approx\bm{f}(\bm{x})+\bm{J}\delta\bm{x}$。

\textbf{例}：$\bm{f}(x,y)=\begin{bmatrix}x^{2}y\\ \sin x\end{bmatrix}$，则
\begin{equation}
\bm{J}=\begin{bmatrix}2xy&x^{2}\\ \cos x&0\end{bmatrix}.
\end{equation}

\subsection{海森矩阵与极值}

\begin{definition}[海森矩阵]
标量函数 $f:\R^{n}\to\R$ 的二阶偏导矩阵
\begin{equation}
\bm{H}=\frac{\partial^{2}f}{\partial\bm{x}^{2}},\qquad H_{ij}=\frac{\partial^{2}f}{\partial x_i\partial x_j}
\end{equation}
叫海森矩阵。
\end{definition}
若 $\nabla f(\bm{x}^{*})=\bm{0}$ 且 $\bm{H}(\bm{x}^{*})$ 正定，则 $\bm{x}^{*}$ 是局部极小点；若负定则是极大点；若不定则是鞍点。这个判别法与一元函数“二阶导数为正”完全对应。

\subsection{最小二乘与正规方程}

\begin{definition}[最小二乘]
给定超定方程组 $\bm{A}\bm{x}\approx\bm{b}$（$\bm{A}\in\R^{m\times n}$，$m>n$），最小二乘解是使残差平方和最小的 $\bm{x}$：
\begin{equation}
\bm{x}^{*}=\argmin_{\bm{x}}\ \norm{\bm{A}\bm{x}-\bm{b}}^{2}.
\end{equation}
\end{definition}

推导：令 $F(\bm{x})=\norm{\bm{A}\bm{x}-\bm{b}}^{2}=(\bm{A}\bm{x}-\bm{b})^{\top}(\bm{A}\bm{x}-\bm{b})$。展开：
\begin{equation}
F(\bm{x})=\bm{x}^{\top}\bm{A}^{\top}\bm{A}\bm{x}-2\bm{b}^{\top}\bm{A}\bm{x}+\bm{b}^{\top}\bm{b}.
\end{equation}
对 $\bm{x}$ 求梯度并令其为零：
\begin{equation}
\nabla F=2\bm{A}^{\top}\bm{A}\bm{x}-2\bm{A}^{\top}\bm{b}=\bm{0},
\end{equation}
得到\textbf{正规方程（normal equation）}
\begin{equation}
\boxed{\bm{A}^{\top}\bm{A}\,\bm{x}=\bm{A}^{\top}\bm{b}}
\label{eq:02-normal}
\end{equation}
若 $\bm{A}^{\top}\bm{A}$ 可逆，则 $\bm{x}=(\bm{A}^{\top}\bm{A})^{-1}\bm{A}^{\top}\bm{b}$。

\textbf{几何解释}：$\bm{A}\bm{x}^{*}$ 是 $\bm{b}$ 在 $\bm{A}$ 的列空间上的正交投影，残差 $\bm{A}\bm{x}^{*}-\bm{b}$ 垂直于列空间。这就是“最小二乘”名字的来源。

\subsection{高斯--牛顿法}

实际问题里残差通常是非线性的：$\bm{r}(\bm{x})$，目标是最小化 $\frac12\norm{\bm{r}(\bm{x})}^{2}$。对残差做一阶泰勒展开：
\begin{equation}
\bm{r}(\bm{x}+\delta\bm{x})\approx\bm{r}(\bm{x})+\bm{J}\delta\bm{x},
\end{equation}
代入目标函数：
\begin{equation}
\frac12\norm{\bm{r}+\bm{J}\delta\bm{x}}^{2}
=\frac12\left(\bm{r}^{\top}\bm{r}+2\bm{r}^{\top}\bm{J}\delta\bm{x}+\delta\bm{x}^{\top}\bm{J}^{\top}\bm{J}\delta\bm{x}\right).
\end{equation}
对 $\delta\bm{x}$ 求导并令其为零：
\begin{equation}
\bm{J}^{\top}\bm{J}\,\delta\bm{x}=-\bm{J}^{\top}\bm{r},
\label{eq:02-gn}
\end{equation}
这就是\textbf{高斯--牛顿方程}。解出 $\delta\bm{x}$ 后更新 $\bm{x}\leftarrow\bm{x}+\delta\bm{x}$，重复直到收敛。它是式~\eqref{eq:02-normal} 的非线性版本，也是第 7 章 ICP、第 13 章因子图、第 15 章 Cartographer 扫描匹配共同的核心。

高斯--牛顿在残差较大或 $\bm{J}^{\top}\bm{J}$ 接近奇异时会不稳定。Levenberg--Marquardt（LM）给它加了一个阻尼项：
\begin{equation}
\left(\bm{J}^{\top}\bm{J}+\lambda\bm{I}\right)\delta\bm{x}=-\bm{J}^{\top}\bm{r},
\label{eq:02-lm}
\end{equation}
$\lambda$ 大时退化为梯度下降（稳定但慢），$\lambda$ 小时退化为高斯--牛顿（快但可能发散）。实际实现中 $\lambda$ 会根据每步误差是变大还是变小自适应调整。

\section{线性代数补充}\label{sec:02-linalg}

\subsection{特征值与特征向量}

\begin{definition}[特征值与特征向量]
对方阵 $\bm{A}\in\R^{n\times n}$，若存在非零向量 $\bm{v}$ 和标量 $\lambda$ 满足
\begin{equation}
\bm{A}\bm{v}=\lambda\bm{v},
\end{equation}
则称 $\lambda$ 为特征值，$\bm{v}$ 为对应的特征向量。
\end{definition}
它的几何意义是：矩阵 $\bm{A}$ 作用在 $\bm{v}$ 上不改变方向，只把长度缩放 $\lambda$ 倍。

对称矩阵（如协方差矩阵、$\bm{J}^{\top}\bm{J}$）有非常好的性质：
\begin{enumerate}
  \item 所有特征值都是实数；
  \item 不同特征值对应的特征向量相互正交；
  \item 可以正交对角化：$\bm{A}=\bm{Q}\bm{\Lambda}\bm{Q}^{\top}$，其中 $\bm{Q}$ 的列是单位正交特征向量，$\bm{\Lambda}=\diag(\lambda_1,\dots,\lambda_n)$。
\end{enumerate}

在定位里，特征值常被用来判断“可观测性”。例如扫描匹配的信息矩阵 $\bm{\Lambda}=\bm{J}^{\top}\bm{J}$，若某个特征值接近零，说明对应方向上“怎么动残差都不变”，即该方向退化（长走廊、大平面就是这样）。第 7 章和第 23 章都会用到这个判据。

\subsection{二次型与正定}

形如 $q(\bm{x})=\bm{x}^{\top}\bm{A}\bm{x}$ 的表达式叫二次型。若对任意 $\bm{x}\neq\bm{0}$ 都有 $q(\bm{x})>0$，称 $\bm{A}$ 正定（$\bm{A}\succ0$），等价于所有特征值大于零。正定二次型的等高线是椭球，椭球的半轴长度是 $1/\sqrt{\lambda_i}$——这就是协方差矩阵“椭圆”的严格解释。

\subsection{奇异值分解}

任意矩阵 $\bm{A}\in\R^{m\times n}$ 可以分解为
\begin{equation}
\bm{A}=\bm{U}\bm{\Sigma}\bm{V}^{\top},
\end{equation}
其中 $\bm{U}\in\R^{m\times m}$、$\bm{V}\in\R^{n\times n}$ 是正交矩阵，$\bm{\Sigma}$ 是对角矩阵（对角元为非负奇异值）。SVD 在点云配准里有关键作用：第 7 章 ICP 的闭式旋转解就是用 $\bm{W}=\bm{U}\bm{\Sigma}\bm{V}^{\top}$ 的 SVD 得到的。

\section{李群与李代数}\label{sec:02-lie}

旋转矩阵有 9 个元素，但只有 3 个自由度（正交约束 6 个）。这种“参数多于自由度”的表示没法直接做“加法更新”——$\bm{R}+\delta\bm{R}$ 一般不再是旋转矩阵。李群/李代数就是解决这个问题的工具。

\subsection{$\SO(3)$ 与 $\so(3)$}

旋转矩阵构成群 $\SO(3)$，它既是群又是光滑流形，所以叫李群。它在单位元处的切空间叫李代数 $\so(3)$，由所有 $3\times3$ 反对称矩阵组成：
\begin{equation}
\so(3)=\left\{\bm{\phi}^{\wedge}\ \middle|\ \bm{\phi}\in\R^{3}\right\},\qquad
\bm{\phi}^{\wedge}=\begin{bmatrix}0&-\phi_3&\phi_2\\ \phi_3&0&-\phi_1\\ -\phi_2&\phi_1&0\end{bmatrix}.
\label{eq:02-hat}
\end{equation}
符号 $(\cdot)^{\wedge}$ 表示“向量变反对称矩阵”，逆运算 $(\cdot)^{\vee}$ 表示“反对称矩阵变回向量”。

指数映射把李代数元素变成旋转矩阵：$\bm{R}=\exp(\bm{\phi}^{\wedge})$。推导它的闭式表达（Rodrigues 公式）。把指数按泰勒级数展开：
\begin{equation}
\exp(\bm{\phi}^{\wedge})=\bm{I}+\bm{\phi}^{\wedge}+\frac{1}{2!}(\bm{\phi}^{\wedge})^{2}+\frac{1}{3!}(\bm{\phi}^{\wedge})^{3}+\cdots
\end{equation}
注意到 $\bm{\phi}^{\wedge}$ 有一个重要性质：若记 $\theta=\norm{\bm{\phi}}$、$\bm{a}=\bm{\phi}/\theta$，则
\begin{equation}
(\bm{a}^{\wedge})^{2}=\bm{a}\bm{a}^{\top}-\bm{I},\qquad (\bm{a}^{\wedge})^{3}=-\bm{a}^{\wedge}.
\end{equation}
第一个式子可以直接验证（对任意单位向量 $\bm{a}$）；第二个由第一个两边乘 $\bm{a}^{\wedge}$ 得到。于是奇数项和偶数项可以分别合并：
\begin{align}
\exp(\theta\bm{a}^{\wedge})
&=\bm{I}+\left(\theta-\frac{\theta^{3}}{3!}+\frac{\theta^{5}}{5!}-\cdots\right)\bm{a}^{\wedge}
+\left(\frac{\theta^{2}}{2!}-\frac{\theta^{4}}{4!}+\cdots\right)(\bm{a}^{\wedge})^{2}\nonumber\\
&=\bm{I}+\sin\theta\,\bm{a}^{\wedge}+(1-\cos\theta)(\bm{a}^{\wedge})^{2}.
\label{eq:02-rodrigues}
\end{align}
这就是\textbf{Rodrigues 公式}。它说明：绕单位轴 $\bm{a}$ 旋转 $\theta$ 角的旋转矩阵只由 $\bm{\phi}=\theta\bm{a}$ 决定。反过来，给定 $\bm{R}$，也能求出 $\bm{\phi}$（对数映射）：
\begin{equation}
\theta=\arccos\left(\frac{\tr(\bm{R})-1}{2}\right),\qquad
\bm{\phi}=\frac{\theta}{2\sin\theta}\left(\bm{R}-\bm{R}^{\top}\right)^{\vee},
\end{equation}
当 $\theta\to0$ 或 $\theta\to\pi$ 时需要特殊处理（代码里通常都有分支）。

\subsection{$\SE(3)$ 与 $\se(3)$}

把平移也纳入进来，$\SE(3)$ 的李代数 $\se(3)$ 是 6 维的：
\begin{equation}
\se(3)=\left\{\bm{\xi}^{\wedge}\ \middle|\ \bm{\xi}=\begin{bmatrix}\bm{\rho}\\ \bm{\phi}\end{bmatrix}\in\R^{6}\right\},\qquad
\bm{\xi}^{\wedge}=\begin{bmatrix}\bm{\phi}^{\wedge}&\bm{\rho}\\\bm{0}^{\top}&0\end{bmatrix}.
\end{equation}
指数映射为
\begin{equation}
\exp(\bm{\xi}^{\wedge})=\begin{bmatrix}\exp(\bm{\phi}^{\wedge})&\bm{J}\bm{\rho}\\\bm{0}^{\top}&1\end{bmatrix},\qquad
\bm{J}=\bm{I}+\frac{1-\cos\theta}{\theta^{2}}\bm{\phi}^{\wedge}+\frac{\theta-\sin\theta}{\theta^{3}}(\bm{\phi}^{\wedge})^{2},
\end{equation}
其中 $\bm{J}$ 叫左雅可比。这里不要求会推导 $\bm{J}$，只要求知道“$\SE(3)$ 上的增量也是 6 维向量”。

\subsection{扰动模型}

在优化中，我们关心的是“给位姿加一个很小的增量”。有两种常见定义：
\begin{align}
\text{左扰动：}&\quad \bm{R}\leftarrow\exp(\delta\bm{\phi}^{\wedge})\bm{R},\\
\text{右扰动：}&\quad \bm{R}\leftarrow\bm{R}\exp(\delta\bm{\phi}^{\wedge}).
\end{align}
两种都对，区别在于雅可比的形式不同。本书约定：位姿用右扰动，即
\begin{equation}
\bm{R}\leftarrow\bm{R}\exp(\delta\bm{\phi}^{\wedge}),\qquad \bm{t}\leftarrow\bm{t}+\bm{R}\,\delta\bm{\rho},
\end{equation}
这样雅可比只和当前位姿有关，推导更整齐。注意 $\delta\bm{\rho}$ 是在机体系下定义的平移增量。

两个旋转复合的对数并不等于两个对数相加，而是有 BCH（Baker--Campbell--Hausdorff）公式：
\begin{equation}
\log\left(\exp(\bm{\phi}_1^{\wedge})\exp(\bm{\phi}_2^{\wedge})\right)
\approx
\begin{cases}
\bm{J}_l(\bm{\phi}_2)^{-1}\bm{\phi}_1+\bm{\phi}_2, & \text{左乘},\\
\bm{J}_r(\bm{\phi}_1)^{-1}\bm{\phi}_2+\bm{\phi}_1, & \text{右乘},
\end{cases}
\end{equation}
当 $\bm{\phi}_2$ 很小时取一阶近似即可。这个公式出现在第 13 章位姿图残差的雅可比推导里。

\subsection{为什么必须在流形上优化}

如果直接对旋转矩阵的 9 个元素做梯度下降，更新后的矩阵不再满足 $\bm{R}^{\top}\bm{R}=\bm{I}$，需要每步重新正交化，既慢又不自然。正确做法是：\textbf{在切空间（李代数）里求增量，再通过指数映射回到流形}：
\begin{equation}
\bm{R}_{k+1}=\bm{R}_k\exp(\delta\bm{\phi}^{\wedge}),
\end{equation}
其中 $\delta\bm{\phi}\in\R^{3}$ 是普通的、可以随便加减的向量。这样优化变量永远是 3 维的，约束自动满足。所有现代 SLAM 后端（g2o、GTSAM、Ceres 的 \lstinline{LocalParameterization}）都是这么做的。

\begin{summary}

\begin{enumerate}
  \item 刚体变换是“旋转 + 平移”，写成 $\SE(3)$ 齐次矩阵后复合就是乘法，逆有闭式表达。
  \item 期望是概率加权平均；方差是偏离期望的平方的期望，有 $\Var(X)=\E[X^{2}]-(\E[X])^{2}$、$\Var(aX+b)=a^{2}\Var(X)$、独立时方差可加三个常用性质。
  \item 协方差矩阵的对角元是方差，非对角元是协方差；它的特征向量给出不确定性的主轴。
  \item 高斯分布由均值与协方差完全确定；马氏距离的平方服从卡方分布，可用于外点剔除。
  \item 贝叶斯公式把“先验 + 似然”变成“后验”，这是所有定位算法的出发点。
  \item 非线性最小二乘 = 泰勒展开 + 正规方程；高斯--牛顿迭代求解，LM 加阻尼保证稳定。
  \item 旋转不能直接相加，要在李代数里求增量再用指数映射更新；Rodrigues 公式给出指数映射的闭式解。
\end{enumerate}

\end{summary}

\begin{problemset}

  \item 掷一枚均匀骰子，求 $X$ 的期望、方差和标准差。若把每个结果乘以 2 再加 1，新的期望和方差是多少？
  \item 证明 $\Cov(X,Y)=\E[XY]-\E[X]\E[Y]$，并说明为什么“协方差为零”不能推出“独立”。
  \item 一维高斯分布中，若 $\mu=2$、$\sigma=3$，求 $P(2\le X\le 5)$（提示：用标准正态分布表，$P(-1\le Z\le1)\approx0.6827$）。
  \item 对 $\bm{f}(\bm{x})=\bm{R}\bm{x}+\bm{t}$，求雅可比 $\partial\bm{f}/\partial\bm{x}$；若 $\bm{R}$ 用增量 $\exp(\delta\bm{\phi}^{\wedge})$ 扰动，求 $\partial\bm{f}/\partial\delta\bm{\phi}$。
  \item 用 Rodrigues 公式计算绕 $z$ 轴转 $90^\circ$ 的旋转矩阵，并与直觉核对。
  \item 设 $\bm{A}=\begin{bmatrix}2&1\\1&2\end{bmatrix}$，求特征值、特征向量，并说明它对应的二次型椭球的半轴方向。
\end{problemset}

\begin{reading}

\begin{itemize}
  \item \cite{barfoot2017state}：第 2--4 章系统讲解概率、估计与李群，适合作为本章的补充教材。
  \item \cite{sola2017eskf}：用 50 多页把四元数、SO(3)、ESKF 讲透，是 IMU 状态估计必读。
  \item \cite{thrun2005probabilistic}：第 2 章讲贝叶斯滤波，第 6 章讲栅格地图，经典且清晰。
  \item 3Blue1Brown 的线性代数系列视频：直观理解特征值、SVD 和矩阵变换。
\end{itemize}


\end{reading}


